% ID: SEL-D0001-OB-E04
% ESTADO: cotejada. Fuente y transcripción revisadas de forma independiente; salida compilada inspeccionada.
% Fuente: Selectividad/sel_2019_matematicas_aplicadas/Titular Junio/Exámenes Mat Apli CCSS A y B.pdf
% SHA-256: fa0cbbddfdf34b1f9fcf8d1090c2d0c3f13e5353e06da22ec354d3345e313584
% Página: 2; opción B; ejercicio 4.
% Transcripción de primera pasada basada en la imagen de la página original.
% Se conservan los nueve datos, su orden y los puntos decimales de la fuente.
% Sin solución ni modificación de datos o condiciones de la fuente.

Los directivos de una empresa desean estimar el tiempo medio que tardan los empleados en llegar al puesto de trabajo desde sus domicilios. Admitimos que dicho tiempo sigue una distribución Normal de desviación típica 8 minutos. Se elige al azar una muestra de 9 empleados de esa empresa, obteniéndose los siguientes resultados, expresados en minutos:
\[
\begin{array}{rrrrrrrrr}
10 & 17 & 8 & 27 & 6 & 9 & 32 & 5 & 21
\end{array}
\]
\begin{enumerate}[label=\alph*),leftmargin=2.4em]
  \item \textbf{(1.5 puntos)} Determine un intervalo de confianza al 92\,\% para la media poblacional.
  \item \textbf{(1 punto)} Con una confianza del 95.5\,\%, \textquestiondown{}qué tamaño muestral mínimo sería necesario para estimar el tiempo medio con un error inferior a 1.5 minutos?
\end{enumerate}
